\section{Odd-Even effect}

\subsection{Zero modes}

\begin{frame}{Trivial solution}
  \(\gamma_m = 0\) is associated with \alert{conserved quantities}.

  \begin{columns}
    \begin{column}{0.5\textwidth}
      \begin{itemize}
        \item \(\gamma_0 = 0\)
              \begin{enumerate}
                \item \(\phi_0 = 1\): particle number
                \item \(\phi_0 = \epsilon(\vv{p})\): energy
              \end{enumerate}
      \end{itemize}
    \end{column}
    \begin{column}{0.51\textwidth}
      \begin{itemize}
        \item \(\gamma_1 = 0\)
              \begin{enumerate}
                \item \(\phi_1 = \e^{\im \theta}\): \(x\)-momentum
                \item \(\phi_1 = \e^{-\im \theta}\): \(y\)-momentum
              \end{enumerate}
      \end{itemize}
    \end{column}
  \end{columns}
\end{frame}

\subsection{Kinematic Constraints}

\begin{frame}
  \begin{columns}
    \begin{column}{0.6\textwidth}
      Near the Fermi surface,
      \begin{equation}
        \epsilon(\vv{p}_1) \approx \epsilon(\vv{p}_2) \approx \epsilon(\vv{p}_1') \approx \epsilon(\vv{p}_2') \approx \epsilon_F.
      \end{equation}
      Thus, \alert{conservation of energy} implies \alert{conservation of momentum magnitude}:
      \begin{equation}
        p_1 + p_2 = p_1' + p_2'.
      \end{equation}
      Three types of collisions that satisfy this constraint:
      \begin{itemize}
        \item Head-on collisions: \(\vv{p}_1 = -\vv{p}_2, \vv{p}_1' = -\vv{p}_2'\)
        \item Forward exchange: \(\vv{p}_1' = \vv{p}_2, \vv{p}_2' = \vv{p}_1\)
        \item No collision: \(\vv{p}_1' = \vv{p}_1, \vv{p}_2' = \vv{p}_2\)
      \end{itemize}
    \end{column}
    \begin{column}{0.4\textwidth}
      \begin{figure}[]
        \centering
        \includegraphics{collision_types.pdf}
        \caption{Relevant collision types}
      \end{figure}
    \end{column}
  \end{columns}
\end{frame}

\begin{frame}
  \begin{description}
    \item [\textbf{No collisions:}] Has no contribution to the collision integral.
    \item [\textbf{Forward exchange \& No collisions:}]
          \begin{itemize}
            \item \(\vv{p}_1 = \vv{p}_1'\), \(\vv{p}_2 = \vv{p}_2'\).
            \item \(\psi(\vv{p}_1) + \psi(\vv{p}_2) - \psi(\vv{p}_1') - \psi(\vv{p}_2') = 0\)\footnote{Recall that this is a multiplier in the collision integral}
            \item \(\therefore\) These collisions has no contribution to the collision integral.
          \end{itemize}
  \end{description}

  Only \alert{head-on collision} contribute to the collision integral.
\end{frame}

\begin{frame}{Head on collision}
  \begin{columns}
    \begin{column}{0.4\textwidth}
      \begin{figure}
        \centering
        \includegraphics{head-on_collision.pdf}
        \caption{Head-on collision angle}
      \end{figure}
    \end{column}
    \begin{column}{0.6\textwidth}
      \begin{equation}
        \theta_2 = \theta_1 + \cpi, \mathand \theta_2' = \theta_1' + \cpi.
      \end{equation}
      Substituting test mode \(m\), and recalling \(\e^{\im m (\theta + \cpi)} = (-1)^m \e^{\im m \theta}\), then
      \begin{multline}
        \psi(\vv{p}_1) + \psi(\vv{p}_2) - \psi(\vv{p}_1') - \psi(\vv{p}_2') \\
        = (1 + (-1)^m)\e^{\im m \theta_1} - (1 + (-1)^m)\e^{\im m \theta_{1'}}.
      \end{multline}
      \vspace{-1.5em}
      \begin{description}
        \item [\(m\) is odd:] the term vanishes; thus the collision integral vanishes.
        \item [\(m\) is even:] the term does not vanish.
      \end{description}
    \end{column}
  \end{columns}
\end{frame}
